\begin{abstract}
In-context learning (ICL) allows a segmentation model to adapt to a new
anatomical structure at inference time from a handful of (image, mask)
context pairs, without any gradient update. Extending ICL segmentation to
full 3D volumes is attractive for clinical CT/MRI workflows but exposes a
resolution--compute trade-off that is more severe than in 2D: dense
cross-attention between a target volume and its contexts scales cubically
with side length, forcing prior work to either downsample the volume,
fall back to per-slice processing, or tile the volume with independent,
context-blind sliding windows. We propose a 3D in-context segmentation
model built around two components. First, a
\emph{bi-axial} image--label attention mechanism lets target and context
tokens exchange information jointly along the spatial-patch axis and the
image/label axis, in contrast to early feature concatenation
(Medverse~\cite{huMedverseUniversalModel2025a}) or one-shot
pixel-shuffle fusion (Iris~\cite{gaoShowSegmentUniversal2025}). Second, a
coarse-to-fine cascade initializes each finer level's label tokens from
the previous level's upsampled prediction and restricts its sliding
window to the previously predicted foreground region, while a small set
of learned register tokens carries hidden state across levels so the
whole cascade can be trained end-to-end with gradients flowing between
resolutions. We evaluate on TotalSegmentator~\cite{wasserthalTotalSegmentatorRobustSegmentation2023}
CT at a fixed 3\,mm spacing, reporting Dice, normalized surface distance,
and compute (wall-clock, FLOPs, peak VRAM) against Medverse and Iris
baselines, and ablate the query prior, the region-restricted fine level,
and cross-level register training. \TODO{fill in headline numbers once
experiments are final: e.g., X\% Dice at Y\% of the compute of the
strongest baseline, and standout categories such as small/thin
structures.} We further test generalization to \TODO{other CT cohorts,
MRI, and far-OOD tasks}, and release code and checkpoints.
\end{abstract}
